\appendix

\section{Proof of Conformal Coverage (Proposition~\ref{prop:coverage})}
\label{app:coverage}

We reproduce the standard split conformal argument for completeness.

\begin{proof}
Let $S_1, \ldots, S_n, S_{\text{test}}$ be iid from distribution $P$ over $\mathbb{R}$.
Define $\hat{q}_\alpha = S_{(\lceil (n+1)(1-\alpha) \rceil)}$, 
the $\lceil (n+1)(1-\alpha) \rceil$-th order statistic of $S_1, \ldots, S_n$.

By exchangeability of $(S_1, \ldots, S_n, S_{\text{test}})$, 
the rank of $S_{\text{test}}$ among $S_1, \ldots, S_n, S_{\text{test}}$ 
is uniformly distributed on $\{1, \ldots, n+1\}$.
Therefore:
\[
  \Pr[S_{\text{test}} > \hat{q}_\alpha] 
  = \Pr\!\left[\mathrm{rank}(S_{\text{test}}) > \lceil(n+1)(1-\alpha)\rceil\right]
  = \frac{\lfloor (n+1)\alpha \rfloor}{n+1}
  \leq \alpha.
\]
Adding the $\frac{1}{n+1}$ slack for the finite-sample correction 
(since the empirical quantile can shift by at most one order statistic) 
gives the stated bound.
\end{proof}

\noindent
In our calibration, the active dataset uses $n=2{,}000$ clean calibration
examples, giving finite-sample slack below $5\times10^{-4}$.

% ─────────────────────────────────────────────────────────────────────────────

\section{TDA Computation Details}
\label{app:tda}

\paragraph{Persistence homology pipeline.}
For each layer $\ell$, the activation tensor $\phi_\ell(x) \in \mathbb{R}^{d_\ell}$ 
is downsampled to a point cloud 
$\mathcal{P} \subset \mathbb{R}^{d_\ell}$ with $|\mathcal{P}| \leq 150$
via deterministic hash-based subsampling.
We then:
\begin{enumerate}[noitemsep]
  \item Construct the Vietoris-Rips filtration on $\mathcal{P}$ 
        with parameter $r \in [0, r_\text{max}]$, $r_\text{max}=50$.
  \item Compute $H_0$ (connected components) and $H_1$ (1-cycles) persistence.
  \item Return the persistence diagram $\mathrm{Dgm}(\mathcal{P})$ 
        as a list of $(b_i, d_i)$ pairs.
\end{enumerate}
We use gudhi v3.12 \cite{gudhi2015} for filtration construction and 
persistent homology computation.

\paragraph{Wasserstein medoid computation.}
Given $N$ diagrams $\{\mathcal{D}_i\}_{i=1}^N$, 
the medoid is:
\[
  i^* = \arg\min_{i} \sum_{j=1}^N W_2(\mathcal{D}_i, \mathcal{D}_j).
\]
We compute exact $W_2$ distances via the auction algorithm 
as implemented in gudhi \cite{gudhi2015}.
For $N=5{,}000$ profile diagrams with $\sim 20$ points each,
full pairwise computation takes $\sim 8$~minutes on a single CPU core.

\paragraph{Expert clustering.}
$K$-medoids clustering with $K=4$ is applied to the $N$ calibration 
persistence diagrams using the Wasserstein PAM algorithm \cite{schubert2021fast}.
Each cluster medoid serves as the routing signature $\mathcal{R}_k^L$
in Eq.~(\ref{eq:routing}).
Expert sub-networks are 3-layer MLPs (Linear $\to$ ReLU $\to$ LayerNorm
$\to$ Linear $\to$ ReLU $\to$ LayerNorm $\to$ Linear) with hidden dimension
$256$, trained on per-attack-cluster activation pools. We use LayerNorm
rather than BatchNorm because the per-cluster pool size at inference can be
as small as 1 (a single L3-rejected input).

% ─────────────────────────────────────────────────────────────────────────────

\section{Campaign Detection Latency}
\label{app:campaign}

To isolate the SACD contribution, we evaluate \emph{detection latency}: 
the number of adversarial queries before L0 mode fires.
We construct a synthetic query stream: 50 clean images, 
then consecutive PGD/FGSM adversarial bursts according to the campaign
scenario definitions emitted by \texttt{run\_campaign\_eval.py}.

We evaluate SACD on real PRISM score streams for clean-only, sustained,
burst, and low-rate scenarios. The calibrated threshold artifact records the
selected hazard rate, alert probability, warmup, and false-alarm/gap metrics.

\paragraph{Calibration requirement.}
BOCPD requires domain-calibrated priors.
With default prior ($\mu_0=0.1$), the detector does not fire within 150
queries because the NIG posterior cannot match the actual canonical clean
score mean ($\sim 0.29$ on the ensemble logit scale; see
\texttt{score\_quantiles.clean.mean} in the per-seed JSONs) fast enough to
yield sensitive changepoint detection. For deployment, $\mu_0$ should be
set to the empirical mean of the calibration-set anomaly scores.

% ─────────────────────────────────────────────────────────────────────────────

\section{Hyperparameter Sensitivity}
\label{app:hyperparams}

\begin{table}[h]
\centering
\small
\begin{tabular}{lccl}
\toprule
Hyperparameter & Default & Range tested & Sensitivity \\
\midrule
$n_\text{cal}$ (calibration size)  & 2000 & 500--3000 & Low ($\pm0.03$ on L3) \\
$n_\text{sub}$ (TDA subsampling)   & 150  & 50--500   & Medium (TPR $\pm4\%$) \\
$K$ (expert clusters)              & 4    & 2--8      & Low (Recovery $\pm2\%$) \\
$\lambda$ (L0 threshold factor)    & 0.8  & 0.6--0.95 & Medium (FPR tradeoff) \\
$p_\text{alert}$ (BOCPD threshold) & 0.50 & 0.30--0.70 & Low (latency $\pm2$ steps) \\
$\delta$ (CUSUM drift floor)       & 0.10 & 0.00--0.50 & High at $\delta=0$ \\
$k_\text{consec}$ (CUSUM consec.)  & 3    & 1--5      & High at $k_\text{consec}=1$ \\
\bottomrule
\end{tabular}
\caption{Hyperparameter sensitivity summary.}
\label{tab:hyper}
\end{table}

% ─────────────────────────────────────────────────────────────────────────────

\section{Appendix-Only: SACD ASR Gap (Demoted Metric)}
\label{app:asr_gap}

The naive ``ASR gap (L0 off vs.\ L0 on)'' metric at $\rho{=}1.0$ is reported
here for completeness; the headline campaign metric is time-to-detect
(\S\ref{sec:campaign}, Table~\ref{tab:campaign}).

\begin{table}[h]
\centering
\small
\begin{tabular}{lccc}
\toprule
Scenario & ASR (L0 off) & ASR (L0 on) & $\Delta$ (pp, paired) \\
\midrule
clean\_only              & $0.000 \pm 0.000$ & $0.000 \pm 0.000$ & $+0.00 \pm 0.00$ \\
sustained\_$\rho{=}0.5$  & $0.015 \pm 0.014$ & $0.011 \pm 0.015$ & $+0.39 \pm 0.88$ \\
sustained\_$\rho{=}0.8$  & $0.012 \pm 0.009$ & $0.009 \pm 0.010$ & $+0.24 \pm 0.54$ \\
sustained\_$\rho{=}1.0$  & $0.011 \pm 0.007$ & $0.009 \pm 0.009$ & $+0.20 \pm 0.45$ \\
burst                    & $0.010 \pm 0.016$ & $0.010 \pm 0.016$ & $+0.00 \pm 0.00$ \\
low\_rate ($\rho{=}0.05$) & $0.004 \pm 0.009$ & $0.000 \pm 0.000$ & $+0.41 \pm 0.91$ \\
\bottomrule
\end{tabular}
\caption{Campaign-stream ASR gap (5-seed mean $\pm$ std). The ceiling is
structurally bounded: the upstream L3 conformal threshold rejects $\sim$99\%
of $\rho{=}1.0$ adversarials within 22 queries, leaving no headroom for an
L0-attributable ASR gap (all gaps $\leq 0.41$~pp). Time-to-detect, not ASR
gap, is the operationally meaningful campaign metric (Table~\ref{tab:campaign}).}
\label{tab:asr_gap}
\end{table}

% ─────────────────────────────────────────────────────────────────────────────

\section{Appendix-Only: Weakened-CW Detector-Fidelity Ablation}
\label{app:weak_cw}

A 5-seed local rerun with weakened CW parameters ($\text{max\_iter}=5$,
$\text{bss}=5$, vs.\ the canonical $\text{max\_iter}=100, \text{bss}=9$
in Table~\ref{tab:main}) shows that the ensemble detector's CW TPR depends
on attack-signature fidelity, not on perturbation magnitude:

\begin{table}[h]
\centering
\small
\begin{tabular}{lccc}
\toprule
CW configuration & Mean L2 distortion & Attack success & Detector TPR \\
\midrule
canonical ($\text{max\_iter}=100$, $\text{bss}=9$) & 0.128 & $1.000$ & $\mathbf{0.938 \pm 0.008}$ \\
weakened ($\text{max\_iter}=5$, $\text{bss}=5$)    & 0.193 & $0.896$ & $0.645 \pm 0.013$ \\
\bottomrule
\end{tabular}
\caption{Stronger CW with smaller distortion is \emph{easier} to detect
than weakened CW with larger distortion. Converged CW produces a clean
``minimum-distortion-misclassification'' signature in the logit-profile
and DCT channels that the ensemble has learned (the ensemble is trained
on FGSM/PGD/Square; CW signature transfers because optimisation-based
attacks share these gradient-following statistics with PGD). Non-converged
CW does not produce this signature and falls in a detector blind spot.}
\label{tab:weak_cw}
\end{table}

% ─────────────────────────────────────────────────────────────────────────────

\section{Appendix-Only: TAMSH Router Diagnosis (F4 Fix)}
\label{app:tamsh_diag}

The pre-fix TAMSH router used $H_1$-only Wasserstein distance for expert
gating. Diagnostic audit (seed 42, $n=200$ PGD inputs) showed:

\begin{enumerate}[leftmargin=*,noitemsep]
  \item $H_1$ medoid diagrams contain only 1--3 points per expert on
        \texttt{layer4} of a CIFAR-native ResNet-18.
  \item 22.7\% of L3-rejected inputs have \emph{empty} $H_1$ diagrams.
  \item The empty-diagram fallback in the Wasserstein implementation
        returns $\sum |b-d|$, which evaluates to the same tiny value for
        all four experts. \texttt{argmin} then defaults to index 0.
  \item Net effect: expert 0 (clean-trained) wins 94\% of routes ($n=163$,
        picks=[154, 0, 8, 1]); routing entropy = 0.336/2.000~bits.
  \item Force-routing the same inputs to each expert in isolation shows
        expert 2 (PGD-trained) achieves 35\% recovery accuracy vs.\
        $\sim$5\% for experts 0/1/3. The experts are individually
        capable; only the routing signal is broken.
\end{enumerate}

\noindent
The F4 fix adds an $H_0$ term to the routing metric (Eq.~\ref{eq:routing}):
$H_0$ medoid diagrams have $\sim$16 points (one per connected component
at large filtration radius) and are sufficiently dense for Wasserstein
gating to be informative. On the 5-seed L3-rejected pool, per-input
oracle (best-of-$K$) ceiling is $39.9\%$ recovery; the combined-metric
router achieves $13.9\%$ ($35\%$ of oracle, $+10.6$~pp over passthrough);
always-route-to-expert-2 achieves $29.3\%$ ($73\%$ of oracle, $+26.0$~pp
over passthrough, clearing the $+15$~pp C4 gate).

% ─────────────────────────────────────────────────────────────────────────────

\section{CIFAR-100 Results}
\label{app:cifar100}

We report CIFAR-100 results following the same protocol as
\S\ref{sec:experiments} (5 seeds, $n=1000$ per seed; CIFAR-native ResNet-18
trained from scratch via \texttt{scripts/pretrain\_cifar\_backbone.py
--dataset cifar100}; conformal calibration on a held-out clean split with
the same three FPR tiers L1/L2/L3 = $0.10, 0.03, 0.005$).

The CIFAR-100 pipeline is launched by \texttt{run\_vastai\_cifar100.sh} and
produces the same set of artifacts: detection (FGSM/PGD/Square/CW/AutoAttack),
adaptive PGD, channel-mask ablation, baselines, campaign streams, and
recovery (\texttt{tamsh}, \texttt{tamsh\_uniform}, \texttt{tamsh\_force}).
Results land under \texttt{models/cifar100/} and \texttt{experiments/*}
with the \texttt{cifar100} tag and feed
\texttt{scripts/build\_paper\_tables.py} which emits combined CIFAR-10 +
CIFAR-100 rows into the same five table files used in the main text.
